\subsection{测度与有界连续函数}

回忆（GT, IX, §1, No.~5, 定义 4），拓扑空间 T 称为完全正则的，是指它可一致化且为豪斯多夫空间。这等价于（同上，命题 3）T 同胚于某个紧空间的子空间。若 T 完全正则，则 T 上每个正的下半连续函数 f 都是由满足 \(\mathcal{C}_{+}^{b}(T)\) 的 \(\leqslant f\) 元素构成的递增有向集的上包络；每个正的有界上半连续函数 g 都是由满足 \(\mathcal{C}_{+}^{b}(T)\) 的 \(\geqslant g\) 元素构成的递减有向集的下包络（同上，§1, No.~6, 命题 5）。我们将需要以下引理：

引理。 --- 设 T 是完全正则空间，K 是 T 的紧子集，U 是包含 K 的 T 的开子集。

\begin{enumerate}
\def\labelenumi{\alph{enumi})}
\item
  存在开子集 \(\mathbf{U}'\)（属于 \(\mathbf{T}\)），使得 \(\mathbf{K} \subset \mathbf{U}' \subset \overline{\mathbf{U}'} \subset \mathbf{U}\)。
\item
  设 \(f\) 是定义在 \(K\) 上、取值于区间 \(I\)（属于 \(\mathbf{R}\)）（分别地，取值于 \(\mathbf{C}\)）的连续函数。存在定义在 \(f'\) 上、取值于 \(T\)（分别地，取值于 \(I\)）的有界连续函数 \(\mathbf{C}\)，它延拓 \(f\) 且在 \(T - U\) 上为零。
\end{enumerate}

只须处理 T 是紧空间 X 的子空间的情形。设 V 是 X 的开子集，满足 \(V \cap T = U\)；记 \(V'\) 为 X 中包含 K 且满足 \(\overline{V'} \subset V\) 的开集，记 g 为 X 上取值于 I（分别地，取值于 C）、延拓 f 且在 X - V 上为零的连续函数（GT, IX, §4, No.~1, 命题 1）。取 \(U' = V' \cap T\) 即满足 a)，取 \(f'\) 为 g 在 T 上的限制即满足 b)。

命题 1. --- 设 T 是完全正则空间。

\begin{enumerate}
\def\labelenumi{\alph{enumi})}
\tightlist
\item
  设 \(\mu\) 是 \(\mathbf{T}\) 上的正测度，\(f\) 是满足 \(\geqslant 0\) 且定义在 \(\mathbf{T}\) 上的数值函数，并且是下半连续的（分别地，是上半连续的、有限的、有紧支集的）。则
\end{enumerate}

\[
\mu^ {\bullet} (f) = \sup _ {g \in I _ {f}} \mu^ {\bullet} (g) \quad (\text { resp. } \mu^ {\bullet} (f) = \inf _ {g \in S _ {f}} \mu (g)),\tag{1}
\]

其中 \(I_{f}\)（分别地，\(S_{f}\)）表示满足 \(0 \leqslant g \leqslant f\)（分别地，\(g \geqslant f\)）的有界连续函数 g 所成的集合。

\begin{enumerate}
\def\labelenumi{\alph{enumi})}
\setcounter{enumi}{1}
\tightlist
\item
  设 \(\theta\) 是 \(\mathbf{T}\) 上的复测度，\(f\) 是满足 \(\geqslant 0\) 且定义在 \(\mathbf{T}\) 上的数值函数，并且是下半连续的。则
\end{enumerate}

\[
| \theta | ^ {\bullet} (f) = \sup _ {g} | \theta (g) |,\tag{2}
\]

其中 \(g\) 遍历有界且对 \(|\theta|\) 可积的连续复函数所成集合，并满足 \(|g| \leqslant f\)。

公式 (1) 中的第一个公式显然成立，因为 \(I_{f}\) 是由连续函数组成的递增有向集，其上包络为 f，可以应用 §1 第 6 号命题 5。如果证明 \(S_{f}\) 含有一个对 \(\mu\) 可积的有界连续函数，同一命题就能推出第二个公式。于是，设 K 是 f 的支集，M 是 f 的上确界；由于 K 紧，M 有限（GT, IV, §6, No.~2, Th. 3）。取包含 K 且满足 \(\mu^{\bullet}(U) < +\infty\) 的开集 U；由（引理）存在取值于 {[}0, M{]}、在 K 上等于 M 且在 U 外为零的连续函数 g；于是 \(g \in S_{f}\)，并且 \(\mu^{\bullet}(g) \leqslant M\mu^{\bullet}(U) < +\infty\)。

转到 \(b\)。显然只须证明 \(|\theta|^{\bullet}(f) \leqslant \sup |\theta(g)|\)。

设 \(a\) 和 \(b\) 是满足 \(a < b < |\theta|^{\bullet}(f)\) 的两个实数。由 (1)，存在函数 \(h \in \mathcal{C}_{+}^{b}(\mathrm{T})\)，使得 \(h \leqslant f\) 且 \(|\theta|^{\bullet}(h) > b\)；记 \(h\) 的上确界为 M。根据 \(|\theta|^{\bullet}\) 的定义（§1, No.~2, 定义 4），存在 T 的紧子集 K，使得 \(|\theta|_{\mathrm{K}}^{\bullet}(h_{\mathrm{K}}) > b\)。于是存在 K 上的连续复函数 \(j\)，使得 \(|j| \leqslant h_{\mathrm{K}}\) 且 \(|\theta_{\mathrm{K}}(j)| > b\)（第 III 章 §1 第 6 号）；取包含 K 且满足 \(|\theta|^{\bullet}(\mathrm{U} - \mathrm{K}) \leqslant \frac{b - a}{\mathrm{M}}\) 的开集 U（§1, No.~9, 命题 13 和 14）；将 \(j\) 延拓为 T 上在 U 外为零的连续复函数 \(k\)（引理）；对每个 \(t \in \mathrm{T}\)，置

\[
g (t) = \left\{ \begin{array}{l l} k (t) & \text { if } | k (t) | \leqslant h (t) \\ \frac {k (t)}{| k (t) |} h (t) & \text { if } | k (t) | > h (t). \end{array} \right.\tag{3}
\]

显然 \(|g| \leqslant h \leqslant f\)，且 \(g = j\) 在 \(\mathbf{K}\) 上成立，因此 \(\left|\theta_{\mathrm{K}}(j) - |\theta(g)|\right| = \left|\left|\theta(j^0)\right| - |\theta(g)|\right| \leqslant |\theta|^{\bullet}(|j^0 - g|) \leqslant \mathbf{M} \cdot |\theta|^{\bullet}(\mathbf{U} - \mathbf{K}) \leqslant b - a\)，所以 \(|\theta(g)| > a\)。另一方面，我们来证明 \(g\) 是连续函数：

由于 \(a\) 只受条件 \(a < |\theta|^{\bullet}(f)\) 的约束，这就说明公式 (2) 的右端 \(\geqslant\) 左端，从而得到命题。现在令 \(F\)（分别地，\(F'\)）为满足 \(t \in T\)、\(|k(t)| \leqslant h(t)\)（分别地，\(|k(t)| \geqslant h(t)\)）的集合。这些集合是闭集，且它们的并为 \(T\)，因此只须证明 \(g_{F}\) 和 \(g_{F'}\) 连续：现在 \(g_{F} = k_{F}\) 显然成立，而 \(g_{F'}\) 在 \(k(t) \neq 0\) 的点也显然连续；另一方面，若 \(t \in F'\) 且 \(k(t) = 0\)，则同样有 \(h(t) = 0\)，而不等式 \(|g| \leqslant h\) 蕴含 \(g\) 在点 \(t\) 连续。

备注。 --- 1) 设 \(f\) 是正的下半连续函数，\(J_f\) 是在某个对 \(\mu\) 可积的开集外为零且被 \(f\) 控制的正有界连续函数所成的集合。可以证明，\(f\) 是 \(J_f\) 的上包络，并且 \(\mu^\bullet(f) = \sup_{g \in L} \mu(g)\)。

\begin{enumerate}
\def\labelenumi{\arabic{enumi})}
\setcounter{enumi}{1}
\tightlist
\item
  若测度 \(\mu\) 有界，则公式 \(\mu^{\bullet}(f) = \inf_{g\in \mathbb{S}_f}\mu (g)\) 显然
\end{enumerate}

对每个上半连续、正且有界的函数 f 都成立。

命题 2. --- 设 \(\eta\) 和 \(\eta'\) 是完全正则空间 \(T\) 上的两个复测度，并且对每个函数 \(\eta(f) = \eta'(f)\)，对 \(f \in \mathcal{C}^b(T)\) 和 \(|\eta|\) 可积，则 \(|\eta'|\)。则 \(\eta = \eta'\)。

令 \(\theta = \eta - \eta'\)，重新采用命题 1 第二部分的证明。可以要求开集 U 对 \(|\eta|\) 和 \(|\eta'|\) 都可积。于是函数 \(g\) 对这两个测度都可积，而关系 \(\theta(g) = 0\) 蕴含 \(a < 0\)；因此 \(|\theta|^{\bullet}(f) = 0\) 对每个正的下半连续函数 \(f\) 成立，最后取 \(|\theta| = 0\) 得到 \(f = +\infty\)。

命题 3. --- 设 \(\mu\) 是完全正则空间 \(T\) 上的正测度，且 \(p \in [1, +\infty[\)。空间 \(\mathcal{H}\) 由支集包含于某个对 \(f \in \mathcal{C}^b(T)\) 可积的开集的函数 \(\mu\) 构成，在 \(\mathcal{L}^p(\mu)\) 中稠密。

由 §1 第 10 号命题 15，只须证明：若 K 是 T 中的紧集，且 g 是 \(\mathcal{C}_{+}(K)\) 中介于 0 和 1 之间的函数延拓到 T 后的函数，则存在函数 \(f \in \mathcal{C}_{+}^{b}(T)\)，其支集包含于某个对 \(\mu\) 可积的开集，并使 \(\|f - g\|_{p}\) 任意小。现在取数 \(\varepsilon\)（\textgreater0）、K 的开邻域 U，使得 \(\mu^{\bullet}(U - K) < \varepsilon\)、K 的开邻域 V，使得 \(\overline{V} \subset U\)，以及取值于 {[}0,1{]}、连续、在 K 上等于 g 且在 V 外为 0 的函数 f（引理）。于是函数 \(|f - g|^{p}\) 被 \(\varphi_{U-K}\) 控制；因此 \(\|f - g\|_{p} \leqslant \varepsilon^{1/p}\)，这就证明了命题。

备注 3)。 --- 对取值于 Banach 空间 F 的函数也有类似结论：\(\mathcal{H} \otimes \mathrm{F}\) 作为 \(\mathcal{C}^b (\mathrm{T};\mathrm{F})\) 的子空间，在 \(\mathcal{L}_{\mathrm{F}}^{p}(\mu)\) 中稠密。

命题 4. --- 要使完全正则空间 \(\theta\) 上的有界复测度 \(T\) 为正，必要且充分的条件是对每个函数 \(\theta(f) \geqslant 0\) 都有 \(f \in \mathcal{C}_{+}^{b}(T)\)。

必要性显然。为证明充分性，重新采用前一命题的证明，取 p = 1 且 \(\mu = |\theta|\)；记号相同，关系 \(\mu^{\bullet}(|f - g|) \leqslant \varepsilon\) 与不等式 \(\theta(f) \geqslant 0\) 蕴含 \(\theta_{\mathrm{K}}(g_{\mathrm{K}}) = \theta(g) \geqslant -\varepsilon\)；由于 \(g_{K}\) 是 \(\mathcal{C}(K)\) 中任意一个介于 0 和 1 之间的元素，测度 \(\theta_{K}\) 为正；K 任意，这意味着 \(\theta\) 为正。

